%% kapa: exam=1 round=25 year=2014 period=2 form=A subject=appraisal_law subject_ko=부동산관계법규 seq=10 number=10
%% source: Q-Net 감정평가사 기출문제 (cst00309) artlSeq=5101207 file=extracted/1008/5101207/2014년 제25회 감정평가사 1차시험 기출문제/2교시(A형).hwp
%% conversion: kordoc 4.24.51 -> tools/gen1.py (auto); review=proofread-sonnet; flags=none
\begin{question}{10}
국토의 계획 및 이용에 관한 법령상 반드시 지구단위계획구역으로 지정해야 하는 지역에 해당하지 않는 것은? (단, 당해 지역에 토지 이용과 건축에 관한 계획이 수립되어 있는 경우가 아님)
\begin{choices}
  \item 개발제한구역ㆍ도시자연공원구역ㆍ시가화조정구역 또는 공원에서 해제되는 구역 중 계획적인 개발 또는 관리가 필요한 지역
  \item 「택지개발촉진법」에 따라 지정된 택지개발지구에서 시행되는 사업이 끝난 후 10년이 지난 지역
  \item 도시지역의 체계적ㆍ계획적인 개발 또는 관리가 필요한 지역으로서 녹지지역에서 주거지역으로 변경되는 지역 중 그 면적이 30만제곱미터 이상인 지역
  \item 도시지역의 체계적ㆍ계획적인 개발 또는 관리가 필요한 지역으로서 녹지지역에서 공업지역으로 변경되는 지역 중 그 면적이 30만제곱미터 이상인 지역
  \item 「도시 및 주거환경정비법」에 따라 지정된 정비구역에서 시행되는 사업이 끝난 후 10년이 지난 지역
\end{choices}
\end{question}
\answer{1}
